% !TEX root = ../Main.tex
\chapter{分式处理}

分式的处理核心是\textbf{"把分子/分母拆成与另一方相关的部分"}——让未知的分式化为"常数 + 简单分式"。常见手段有：分离常数、增减项凑分母、分子裂成分母相关的组合。本章以两个经典场景展开：有理分式估阶、含参分式的不等式证明。

\section{分离常数}

\state{分离常数的动机}{
    对于形如 $\dfrac{P(x)}{Q(x)}$ 的分式（$\deg P \ge \deg Q$），通过多项式除法把它写成
    \[
    \frac{P(x)}{Q(x)} = R(x) + \frac{r(x)}{Q(x)},\qquad \deg r < \deg Q.
    \]
    这样\textbf{分子分母的次数差异被显式分离出来}：$R(x)$ 是主阶部分、余项 $\dfrac{r(x)}{Q(x)}$ 是"次要"扰动，便于分析单调性、估算上下界。
}

\ex[分离常数估阶]{
    已知 $a_n = \dfrac{3 n^2 - n - 2}{2 n^2 + 2 n}$（$n \in \mathbb N^+$）。求证：$a_n > n - 1 + \sum\limits_{i=2}^{n} \ln\!\dfrac{2}{i^2 - 1}$ 对所有 $n \ge 2$ 成立。

    \textit{（提示：利用 $\ln(1+x) < x$（$x > 0$）。）}
}

\sol{
    \textbf{第一步：分离常数。} 对 $a_n$ 作多项式长除：分子 $3 n^2 - n - 2 = \dfrac{3}{2}(2 n^2 + 2 n) - 4 n - 2$，故
    \[
    a_n = \frac{3}{2} + \frac{-4 n - 2}{2 n^2 + 2 n} = \frac{3}{2} - \frac{4 n + 2}{2 n(n+1)} = \frac{3}{2} - \frac{2 n + 1}{n(n+1)}.
    \]

    进一步\textbf{裂项}：$\dfrac{2 n + 1}{n(n+1)} = \dfrac{1}{n} + \dfrac{1}{n+1}$，故
    \[
    a_n = \frac{3}{2} - \frac{1}{n} - \frac{1}{n + 1}.
    \]

    \textbf{第二步：化简右端的 $\ln$ 求和。}
    \[
    \frac{2}{i^2 - 1} = \frac{2}{(i-1)(i+1)} = \frac{1}{i-1} - \frac{1}{i+1}\quad\text{（裂项）}.
    \]
    这暗示 $\dfrac{2}{i^2 - 1}$ 可表示为相邻两项之差；但对数求和里这个裂项不直接消，需要另一思路——\textbf{用 $\ln(1+x) < x$ 估计}：
    \[
    \ln \frac{2}{i^2 - 1} = \ln\!\left(1 + \frac{2 - (i^2 - 1)}{i^2 - 1}\right) = \ln\!\left(1 - \frac{i^2 - 3}{i^2 - 1}\right).
    \]
    当 $i \ge 2$ 时 $\dfrac{2}{i^2 - 1} \le \dfrac{2}{3} < 1$，即 $\ln \dfrac{2}{i^2 - 1} < 0$。因此所证 $a_n > n - 1 + \sum \ln \dfrac{2}{i^2 - 1}$ 等价于
    \[
    a_n - (n - 1) > \sum_{i=2}^{n} \ln \frac{2}{i^2 - 1}.
    \]

    \textbf{第三步：左端。} 由第一步
    \[
    a_n - (n - 1) = \frac{3}{2} - \frac{1}{n} - \frac{1}{n+1} - (n-1) = \frac{5}{2} - n - \frac{1}{n} - \frac{1}{n+1}.
    \]

    \textbf{第四步：右端用 $\ln x < x - 1$（即 $\ln(1+y) < y$，取 $y = x - 1$）。}
    \[
    \sum_{i=2}^{n} \ln \frac{2}{i^2 - 1} < \sum_{i=2}^{n}\!\left(\frac{2}{i^2 - 1} - 1\right) = \sum_{i=2}^{n} \frac{2}{i^2 - 1} - (n - 1).
    \]
    而 $\sum_{i=2}^{n} \dfrac{2}{i^2 - 1} = \sum_{i=2}^{n}\!\left(\dfrac{1}{i-1} - \dfrac{1}{i+1}\right) = 1 + \dfrac{1}{2} - \dfrac{1}{n} - \dfrac{1}{n+1} = \dfrac{3}{2} - \dfrac{1}{n} - \dfrac{1}{n+1}$（裂项相消）。故
    \[
    \sum_{i=2}^{n} \ln \frac{2}{i^2 - 1} < \frac{3}{2} - \frac{1}{n} - \frac{1}{n+1} - (n - 1) = \frac{5}{2} - n - \frac{1}{n} - \frac{1}{n+1} = a_n - (n - 1).
    \]

    即 $a_n - (n - 1) > \sum_{i=2}^{n} \ln \dfrac{2}{i^2 - 1}$，移项即得证。$\blacksquare$
}

\rmkb{
    本例综合了多种技巧：\textbf{分式的分离常数 + 裂项 + 对数的放缩估计 ($\ln(1+x) < x$)}。分离常数一步把复杂有理式化为"$\dfrac{3}{2}$ + 两个简单裂项分式"，为后续比较估算铺平了道路。没有这一步，式子难以下手。
}

\section{增减项配凑}

\state{增减同一项使分子向分母靠拢}{
    对于 $\dfrac{f(n)}{g(n)}$，若 $f(n)$ 与 $g(n)$ 结构接近（如相同次数），常用\textbf{"加一项、减一项"}的技巧：
    \[
    \frac{f(n)}{g(n)} = \frac{g(n) + [f(n) - g(n)]}{g(n)} = 1 + \frac{f(n) - g(n)}{g(n)}.
    \]
    特别是在\textbf{极限 / 估阶}问题中，这把"两个相近大量之比"变成"$1$ 加一个小量"，更容易分析。
}

\ex[配凑出 $1$]{
    设 $n \to \infty$。求 $\dfrac{n^2 + 3 n + 1}{n^2 + 1}$ 的极限。
}

\sol{
    \[
    \frac{n^2 + 3 n + 1}{n^2 + 1} = \frac{(n^2 + 1) + 3 n}{n^2 + 1} = 1 + \frac{3 n}{n^2 + 1} \xrightarrow{n \to \infty} 1 + 0 = 1.
    \]
    （把 $n^2 + 1$ 配出来，剩下的 $\dfrac{3 n}{n^2 + 1}$ 分母高一次，$\to 0$。）
}

\rmkb{
    \textbf{分式处理的通用原则}：让式子\textbf{更"像"已知的结构}。
    \begin{itemize}
        \item \textbf{分离常数}—— 把 $\deg P \ge \deg Q$ 的分式拆成整式 + 真分式；
        \item \textbf{增减项配凑}—— 让分子出现 $g(n)$，得到 $1 + \dfrac{\text{余项}}{g(n)}$；
        \item \textbf{裂项}—— 把分式拆成两个更简单分式之差（本章例 11）。
    \end{itemize}
    这三招的共同精神：\textbf{把未知的分式改写为"常数 + 简单分式"的形式，方便分析}。
}
